\begin{definition}[Сильная связность]
Пишем \(u\leadsto v\), если существует направленный маршрут из \(u\)
в \(v\), включая маршрут длины ноль. Вершины сильно связаны,
\(u\sim v\), если \(u\leadsto v\) и \(v\leadsto u\).
Орграф сильно связен, если все его вершины попарно сильно связаны.
\end{definition}
\begin{theorem}
Сильная связность — отношение эквивалентности. Его классы называются
компонентами сильной связности (КСС); каждый класс индуцирует сильно связный подграф.
\end{theorem}
\begin{proof}
Рефлексивность следует из нулевого маршрута; симметричность заложена
в определении. Если \(u\sim v\) и \(v\sim w\), склейка маршрутов даёт
\(u\leadsto w\) и \(w\leadsto u\), то есть транзитивность.
Если \(u,v\) лежат в одной КСС, то каждая вершина \(x\) маршрута
\(u\leadsto v\) достижима из \(u\) и достигает \(u\) через \(v\).
Значит, \(x\) принадлежит той же КСС и маршрут не выходит из класса.
\end{proof}
\begin{definition}[Слабая связность]
Орграф слабо связен, если связен его неориентированный граф, полученный
забыванием направлений. Сильная связность влечёт слабую, обратное неверно:
одна дуга на двух вершинах даёт контрпример.
Ациклический орграф не имеет направленных циклов положительной длины.
\end{definition}
\begin{theorem}[Топологическая сортировка]
Орграф ацикличен тогда и только тогда, когда существует биекция
\(\varphi:V\to\{1,\ldots,n\}\), для которой
\(uv\in E\Rightarrow\varphi(u)<\varphi(v)\).
\end{theorem}
\begin{proof}
В непустом ациклическом орграфе есть вершина без входящих дуг.
Иначе, выбирая предшественника, затем его предшественника и так далее,
в конечном множестве получим повторение вершины и направленный цикл.
Ставим вершину без входящих дуг первой, удаляем её и повторяем в оставшемся
ациклическом орграфе. При удалении вершины все её предшественники уже
стоят в списке, значит, любая дуга идёт вперёд.
Обратно, вдоль направленного цикла значения \(\varphi\) должны были бы
строго возрастать и вернуться к исходному значению, что невозможно.
\end{proof}
\begin{cor}[Достижимость как порядок]
В ациклическом орграфе отношение \(u\preceq v\), заданное существованием
пути из \(u\) в \(v\), включая нулевой, является частичным порядком.
\end{cor}
\begin{proof}
Рефлексивность даёт нулевой путь, транзитивность — склейка путей.
Взаимная достижимость разных вершин давала бы направленный цикл,
поэтому отношение антисимметрично.
\end{proof}
\textbf{Дополнение по списку вопросов.} В скане следующие определения
не выписаны; они дополняют лекционный порядок на ациклическом орграфе.
\begin{definition}[Конденсация и порядок]
Конденсация \(D^*\) имеет вершинами КСС; дуга \(C\to B\), \(C\ne B\),
проводится, если в исходном орграфе есть дуга из \(C\) в \(B\).
На КСС зададим \(C\preceq B\), если \(C=B\) либо в конденсации есть
направленный путь из \(C\) в \(B\).
\end{definition}
\begin{theorem}
Конденсация ациклична, а \(\preceq\) является частичным порядком.
\end{theorem}
\begin{proof}
Если различные КСС образуют ориентированный цикл, то, используя
сильную связность внутри классов и дуги между ними, из любой вершины
их объединения можно попасть в любую другую. Все они должны составлять
один класс — противоречие. Рефлексивность и транзитивность порядка
следуют из достижимости. Если \(C\preceq B\preceq C\), то при \(C\ne B\)
получается замкнутый направленный маршрут и, после удаления повторений,
цикл в конденсации. Значит, \(C=B\), что даёт антисимметричность.
\end{proof}
\begin{definition}[Диаграмма порядка]
Диаграмма Хассе содержит дугу \(C\to B\), если \(C\prec B\), но нет
класса \(F\), для которого \(C\prec F\prec B\).
Это отношение покрытия, а не все сравнимые пары.
\end{definition}
\begin{lemma}
Диаграмма Хассе ациклична и имеет ту же достижимость, что и конденсация.
\end{lemma}
\begin{proof}
Цикл противоречит строгому возрастанию в частичном порядке.
Для \(C\prec B\) выберем цепочку между ними максимальной длины;
конечность множества гарантирует её существование. Соседние элементы
цепочки образуют покрытие, иначе можно вставить промежуточный элемент.
Поэтому сравнимость реализуется путём по диаграмме. Обратное следует из транзитивности.
\end{proof}
\begin{center}
\begin{tikzpicture}[>=Stealth,every node/.style={circle,draw,minimum size=6mm}]
\node (a) at (0,0) {$A$};\node (b) at (2,0) {$B$};\node (c) at (4,0) {$C$};
\draw[->] (a)--(b);\draw[->] (b)--(c);\draw[->,bend left=35] (a) to (c);
\end{tikzpicture}
\end{center}
Здесь конденсация может содержать все три дуги; в диаграмме Хассе
остаётся только \(A\to B\to C\). Конденсация и диаграмма Хассе —
разные графы; оба ацикличны, но первая может содержать транзитивные дуги.
